% ==================================================
% CHAPTER 3: SYSTEM DESIGN
% ==================================================

\chapter{System Design}
\label{ch:design}

\section{Architecture Overview}

MemoryLane has three client-facing parts and one server
(Figure~\ref{fig:architecture}). The \textbf{Chrome extension} captures
pages; the \textbf{web app} is where users search and manage their library;
and the \textbf{FastAPI backend} processes pages, runs searches and talks to
\textbf{Apache Solr}, which holds both the text index and the vectors. All
browser traffic enters through \textbf{Caddy}, which terminates HTTPS,
serves the built web app and proxies \texttt{/api/*} to the backend.

\begin{figure}[H]
\centering
\begin{tikzpicture}
  % ---- browser -----------------------------------------------------------
  \node[comp, minimum width=30mm] (ext) at (0, 1.1) {Chrome extension\\\scriptsize popup, service worker,\\\scriptsize content scripts};
  \node[comp, minimum width=30mm] (web) at (0,-1.1) {Web app (Svelte)\\\scriptsize search, library,\\\scriptsize account menu};
  \begin{scope}[on background layer]
    \node[group, fill=blue!3, fit=(ext)(web), inner sep=7pt] (browser) {};
  \end{scope}
  \node[font=\small\bfseries, anchor=south] at (browser.north) {User's browser};

  % ---- server ------------------------------------------------------------
  \node[comp, minimum width=22mm, minimum height=36mm] (caddy) at (3.95, 0) {Caddy\\[1mm]\scriptsize HTTPS\\\scriptsize web app\\\scriptsize extension zip\\\scriptsize \texttt{/api} proxy};
  \node[comp, minimum width=32mm, minimum height=36mm] (api) at (7.15, 0)
    {\textbf{FastAPI backend}\\[1mm]\scriptsize auth and accounts\\\scriptsize ingest pipeline\\\scriptsize hybrid retrieval\\\scriptsize snippets, summaries\\\scriptsize problem reports};
  \node[comp, minimum width=29mm] (model) at (11.15, 1.7) {\texttt{bge-small-en-v1.5}\\\scriptsize embeddings, in-process};
  \node[db, minimum width=29mm] (solr)  at (11.15, 0)   {Apache Solr 9.7\\\scriptsize BM25 + HNSW vectors};
  \node[db, minimum width=29mm] (sqlite) at (11.15,-1.7) {SQLite\\\scriptsize accounts, reports};
  \begin{scope}[on background layer]
    \node[group, fill=gray!3, fit=(caddy)(api)(model)(sqlite), inner sep=7pt] (server) {};
  \end{scope}
  \node[font=\small\bfseries, anchor=south] at (server.north) {AWS EC2 $\cdot$ Docker Compose};

  % ---- external services ---------------------------------------------------
  \node[ext] (yt)      at (-0.55,-3.75) {YouTube};
  \node[ext] (cognito) at (4.35,  -3.75) {Amazon Cognito};
  \node[ext] (gemini)  at (7.65,  -3.75) {Google Gemini};
  \node[ext] (ses)     at (10.75, -3.75) {Amazon SES};

  % ---- flows ---------------------------------------------------------------
  \draw[flow] (ext.east) -- node[lbl, above] {HTTPS} (caddy.west |- ext.east);
  \draw[flow] (web.east) -- node[lbl, above] {HTTPS} (caddy.west |- web.east);
  \draw[flow] (caddy) -- (api);
  \draw[flow, <->] (api.east |- model) -- (model.west);
  \draw[flow, <->] (api.east |- solr) -- (solr.west);
  \draw[flow, <->] (api.east |- sqlite) -- (sqlite.west);
  \draw[flow, dashed] (yt.north |- browser.south) -- node[lbl, left] {captions} (yt.north);
  \draw[flow, dashed] (cognito.north |- server.south) -- node[lbl, right] {tokens, email} (cognito.north);
  \draw[flow, dashed] (gemini.north |- server.south) -- node[lbl, right] {summaries} (gemini.north);
  \draw[flow, dashed] (ses.north |- server.south) -- node[lbl, right] {reports} (ses.north);
  \draw[flow, dashed] (1.2,-1.98) |- node[lbl, pos=0.3, right] {sign-in (PKCE)} (cognito.west);
\end{tikzpicture}

\caption{Architecture of the deployed system. Dashed boxes are external services.}
\label{fig:architecture}
\end{figure}

The backend is a single process (one Uvicorn worker) that loads the
embedding model once and shares it between requests. Solr is the only
store for pages; two small SQLite databases hold sign-in links and problem
reports. The design keeps the number of moving parts small enough to run on
one server, a decision recorded in ADR~0004.

\section{Data Flow}

\subsection{Context Diagram (Level 0)}

\begin{figure}[H]
\centering
\begin{tikzpicture}
  \node[proc, minimum size=28mm] (ml) at (0,0) {\textbf{0}\\[0.5mm]MemoryLane};
  \node[entity] (user) at (-5.4, 0)   {Reader};
  \node[entity] (cog)  at (5.0, 2.0)  {Amazon Cognito};
  \node[entity] (gem)  at (5.0, 0)    {Google Gemini};
  \node[entity] (ses)  at (5.0,-2.3)  {Maintainer\\\scriptsize (email via SES)};
  \node[entity] (yt)   at (0,  -3.3)  {YouTube};

  \draw[flow] ([yshift=3mm]user.east) -- node[lbl, above] {pages, queries, deletes} ([yshift=3mm]ml.west);
  \draw[flow] ([yshift=-3mm]ml.west) -- node[lbl, below] {results, snippets, summaries} ([yshift=-3mm]user.east);
  \draw[flow, <->] (ml) -- node[lbl, above, sloped, pos=0.45] {tokens, verified email} (cog.west);
  \draw[flow, <->] (ml) -- node[lbl, above, pos=0.45] {page text, summary} (gem.west);
  \draw[flow] (ml) -- node[lbl, below, sloped, pos=0.45] {problem reports} (ses.west);
  \draw[flow] (yt) -- node[lbl, left] {captions, via the user's browser} (ml);
\end{tikzpicture}

\caption{Context data-flow diagram}
\label{fig:dfd0}
\end{figure}

\subsection{Level 1}

\begin{figure}[H]
\centering
\begin{tikzpicture}[scale=0.8, transform shape, proc/.append style={minimum size=18mm, font=\scriptsize}]
  \newcommand{\dstore}[2]{\makebox[6.6mm][l]{#1}#2}
  \tikzset{wl/.style={lbl, fill=white, inner sep=1pt}}
  % external entities
  \node[entity] (user) at (-5.7, 0)   {Reader};
  \node[entity] (cog)  at (-5.7, 4.2) {Cognito};
  \node[entity] (gem)  at (8.9, 0)    {Gemini};
  \node[entity] (ses)  at (8.6,-4.2)  {SES};

  % processes
  \node[proc] (p1) at (-2.3, 4.2) {\textbf{1}\\Authenticate};
  \node[proc] (p2) at (-2.3, 2.1) {\textbf{2}\\Capture,\\save};
  \node[proc] (p3) at (-2.3, 0.0) {\textbf{3}\\Search};
  \node[proc] (p5) at (-2.3,-2.1) {\textbf{5}\\Manage\\library};
  \node[proc] (p6) at (-2.3,-4.2) {\textbf{6}\\Report a\\problem};
  \node[proc] (p4) at (6.1, 0)    {\textbf{4}\\Summarise};

  % data stores
  \node[store, minimum width=31mm] (d2) at (2.6, 4.2)  {\dstore{D2}{accounts.db}};
  \node[store, minimum width=31mm] (d1) at (2.6, 1.4)  {\dstore{D1}{Solr index}};
  \node[store, minimum width=31mm] (d4) at (2.6,-1.4)  {\dstore{D4}{device store}};
  \node[store, minimum width=31mm] (d3) at (2.6,-4.2)  {\dstore{D3}{reports.db}};

  % flows from the reader
  \foreach \p/\t in {p1/sign-in, p2/page, p3/query, p5/delete, p6/message}
    \draw[flow] (user.east) -- node[lbl, sloped, above, pos=0.5] {\t} (\p.west);

  \draw[flow, <->] (p1) -- node[lbl, above] {code, tokens} (cog);
  \draw[flow, <->] (p1) -- node[lbl, above] {library id} (d2);
  % capture writes to the device store (default) or to Solr
  \draw[flow] (p2.east) -- node[wl, pos=0.3] {chunks (server)} ([yshift=1.5mm]d1.west);
  \draw[flow] (p2.south east) -- node[wl, pos=0.78] {chunks (default)} ([yshift=1.5mm]d4.west);
  % search reads both stores
  \draw[flow] ([yshift=-1.5mm]d1.west) -- node[wl, pos=0.25] {ranked chunks} (p3.north east);
  \draw[flow] ([yshift=-1.5mm]d4.west) -- node[wl, pos=0.25] {ranked chunks} (p3.south east);
  % manage lists and deletes in both
  \draw[flow, <->] (p5.north east) -- node[wl, pos=0.75] {list, delete} ([xshift=6mm]d1.south west);
  \draw[flow, <->] (p5.east) -- node[wl, pos=0.45] {list, delete} ([xshift=6mm]d4.south west);
  % overflow and summaries
  \draw[flow] ([xshift=-5mm]d4.north east) -- node[lbl, left] {oldest past the limit} ([xshift=-5mm]d1.south east);
  \draw[flow] (d1.east) -- node[lbl, sloped, above] {page text} (p4);
  \draw[flow] (d4.east) -- node[lbl, sloped, below] {page text} (p4);
  \draw[flow, <->] (p4) -- node[lbl, above] {summary} (gem);
  \draw[flow] (p6) -- node[lbl, above] {report} (d3);
  \draw[flow] (d3) -- node[lbl, above] {email} (ses);
\end{tikzpicture}

\caption{Level-1 data-flow diagram}
\label{fig:dfd1}
\end{figure}

Every flow that touches D1 carries the verified library id of the caller;
there is no flow by which a client can name another user's id.

\section{Data Model}

\subsection{Pages and Chunks in Solr}

A saved page is stored as $n$ \emph{chunk documents} in one Solr core
(\texttt{memorylane\_v2}). Page-level fields are repeated on every chunk so
that a search hit needs no second lookup; chunk~0 additionally carries the
page vector. Table~\ref{tab:schema} lists the fields.

\begin{table}[H]
\centering
\small
\begin{tabularx}{\textwidth}{>{\ttfamily}p{0.25\textwidth} p{0.2\textwidth} L}
\toprule
\textrm{\textbf{Field}} & \textbf{Type} & \textbf{Purpose} \\
\midrule
id & string & \texttt{\{page\_id\}::\{chunk\_id\}}, so re-saving overwrites in place \\
user\_id & string & Owner's library id; on every query as a filter \\
page\_id & string & \texttt{sha256(user\_id|normalised URL)[:32]} \\
url, domain & string & Normalised address and host \\
title & text\_en & Page title (weighted $\times 2$ in BM25) \\
text & text\_en & The chunk's passage (stemmed, stop-words removed) \\
description & stored & Meta description or first sentences; snippet fallback \\
chunk\_id, total\_chunks & int & Position of the chunk, chunks in the page \\
timestamp & date & When the page was saved \\
content\_hash & string & SHA-256 of the cleaned text; skips unchanged re-saves \\
embedding\_model & string & Model that produced the vectors; a change forces re-embedding \\
embedding & vector(384) & Chunk vector, cosine, HNSW index \\
page\_embedding & vector(384) & Mean of the chunk vectors (chunk~0 only) \\
\bottomrule
\end{tabularx}
\caption{Solr schema (\texttt{schema v2})}
\label{tab:schema}
\end{table}

Automatic field creation is turned off, so a field must be added by a
script before code that writes it is deployed.

\subsection{Relational Data}

Two SQLite databases live in a Docker volume next to Solr's data:

\begin{itemize}
  \item \texttt{accounts.db}: which Cognito sign-ins (\texttt{sub}) share a
  library, the verified email that links them, when each was last checked
  with Cognito, and libraries waiting to be merged after a late
  verification.
  \item \texttt{reports.db}: problem reports and their email delivery
  state, so that a report is never lost when email is temporarily failing
  (ADR~0007).
\end{itemize}

\section{Interface Design}

\subsection{REST API}

\begin{table}[H]
\centering
\small
\begin{tabularx}{\textwidth}{>{\ttfamily}p{0.11\textwidth} >{\ttfamily}p{0.27\textwidth} L r}
\toprule
\textrm{\textbf{Method}} & \textrm{\textbf{Endpoint}} & \textbf{Purpose} & \textbf{Limit} \\
\midrule
GET & /health & Public status of Solr, model, Gemini & -- \\
POST & /memory & Save or update a page & 60/min \\
GET & /memory/status & Is this URL saved? & 120/min \\
POST & /memory/existing & Which of these pages still exist? & 120/min \\
DELETE & /memory/\{id\} & Delete one page & 60/min \\
GET & /memory/\{id\}/related & Pages closest in meaning & 120/min \\
GET & /memories & Library, newest first & 60/min \\
POST & /search & Hybrid search & 60/min \\
POST & /summarize & Summary of one page & 20/min \\
DELETE & /me/data & Delete all the caller's data & 5/min \\
POST & /reports & Report a problem & 5/hour \\
\bottomrule
\end{tabularx}
\caption{API endpoints. All except \texttt{/health} need a Cognito access token.}
\label{tab:api}
\end{table}

\subsection{User Interface}

The web app is deliberately quiet: a dark theme, one search box, and results
that show the title, site, save time and a highlighted snippet. Every action
has a keyboard path (\texttt{/} to search, arrows to move, \texttt{S} to
summarise, \texttt{Shift+Del} to delete with undo). The extension popup
shows whether the current page is saved, one primary button, and the day's
automatic saves with one-click removal.

\section{Design Decisions}

Table~\ref{tab:adrs} lists the architecture decision records kept with the
code. Each record states the alternatives that were rejected and the cost of
the choice.

\begin{table}[H]
\centering
\small
\begin{tabularx}{\textwidth}{>{\ttfamily}p{0.06\textwidth} p{0.36\textwidth} L}
\toprule
\textrm{\textbf{\#}} & \textbf{Decision} & \textbf{Main reason} \\
\midrule
0001 & One Solr core for keyword and vector search & One system for BM25, KNN and filters; no second database to keep in sync \\
0002 & Fuse rankings with Reciprocal Rank Fusion & Rank-based, so no score normalisation and no outlier effects \\
0003 & Drop weak matches with a relevance floor & KNN always returns neighbours, however unrelated \\
0004 & Run everything on one EC2 instance & Cost and simplicity at this scale \\
0005 & Link sign-ins by verified email & One library per person, without trusting the browser \\
0006 & Fetch YouTube captions in the browser & YouTube blocks requests from servers \\
0007 & Store problem reports before emailing & A report survives an email outage \\
0008 & Keep summaries on demand and optional & Search never waits on an external model \\
0009 & Remove the library map & Costly, fragile and rarely useful \\
0010 & Remove ``Export my data'' & Little use; deleting data is the stronger right \\
0011 & Auto-capture keeps only pages that read like content & Rules on content and reading, not an endless site list \\
\bottomrule
\end{tabularx}
\caption{Architecture decision records}
\label{tab:adrs}
\end{table}
